% Jean-Michel Sarlat - 1999

\documentclass[12pt,a4paper]{article}
\usepackage[T1]{fontenc}
\usepackage[utf8]{inputenc}
\usepackage[]{palatino,geometry,french}
\usepackage[dvips]{color}
%=====================================================================
% Définition de \exo ... \finexo

\newcount\exonumber\exonumber=0
\def\Exo{\goodbreak\global\advance\exonumber by 1\par\vskip 15pt
\noindent{\framebox{\bf \color{blue}Exercice \the\exonumber}}}
\def\exo{\begin{center}\Exo\end{center}}
\def\finexo{}
%=====================================================================
\geometry{margin=2cm,head=0.5cm,headsep=10pt,foot=1cm}
\newcommand{\N}{\mathbf{N}}
\parindent0pt
\pagestyle{empty}
\usepackage[np]{numprint}
\begin{document}

\begin{center}
	\sf \Large \color{red}
	Concours général de mathématiques\\
	1999\\
	\rule{0.6\linewidth}{0.3mm}

\end{center}

\exo

Qu'est le volume maximum d'un cylindre, ayant même

axe de révolution qu'un cône donné et intérieur à ce cône~?\\

Quel est le volume maximum d'une boule, centrée sur cet axe et intérieure

au cône ?\\

Comparer les deux maximuns trouvés.

\finexo

\exo

Résoudre dans $\N$ l'équation en $n$ : 

$$(n+3)^n=\sum_{k=3}^{n+2} k^n$$

\finexo

\exo 

Pour quels triangles aux angles tous aigus le quotient du plus

petit coté par le rayon du cercle inscrit est-il maximum~?

\finexo 

\exo

Sur une table trônent \np{1999} bonbons rouges et
\np{6661} bonbons jaunes rendus indiscernables par des emballages uniformes.

Un gourmand applique jusqu'à épuisement du stock l'algorithme ci-dessous~:
\begin{itemize}

	\item[a)] S'il reste des bonbons, il en tire un au hasard, note sa

	couleur, le mange et va en b)~;

	\item[b)] S'il reste des bonbons, il en tire un au hasard et note sa

	couleur~:

	\begin{itemize}

		\item si elle est la même que celle du dernier bonbon avalé, il le

		mange et retourne en~b),

		\item sinon, il le remmaillote,le pose et retourne en a).

	\end{itemize}

\end{itemize}

Montrer que tous les bonbons seront mangés et donner la probabilité pour

que le dernier bonbon mangé soit rouge.

\finexo

\exo

Montrer que les symétriques de chaque sommet d'un triangle par

rapport au coté opposé sont alignés si, et seulement si,

la distance de l'orthocentre au centre du cercle circonscrit est égale à

son diamètre.

\finexo

\end{document}